\documentclass[12pt]{article}
\usepackage[utf8]{inputenc}
\usepackage{amsmath}
\usepackage{graphicx}
\usepackage{booktabs}
\usepackage{float}
\usepackage{geometry}
\geometry{margin=2.5cm}

\title{Avaliação do Modelo KNN para Detecção de Áreas Queimadas}
\author{}
\date{}

\begin{document}

\maketitle
\section{Treinamento e Avaliação do Modelo KNN para Detecção de Áreas Queimadas}

O primeiro modelo avaliado neste trabalho foi o algoritmo \textbf{K-Nearest Neighbors (KNN)}, um método supervisionado amplamente utilizado em tarefas de classificação.

O KNN realiza a classificação com base na similaridade entre amostras. Para cada novo pixel analisado, o algoritmo identifica os $K$ vizinhos mais próximos no conjunto de treinamento e atribui a classe predominante entre eles.

Neste experimento foram utilizados:

\begin{itemize}
    \item $K = 5$ vizinhos;
    \item Padronização das variáveis utilizando \texttt{StandardScaler};
    \item Classificação binária:
    \begin{itemize}
        \item Classe 0: Não queimada;
        \item Classe 1: Queimada.
    \end{itemize}
\end{itemize}

\subsection{Pré-processamento e geração das amostras}

Inicialmente, a abordagem utilizava todos os pixels das imagens como amostras de treinamento. Entretanto, visando reduzir ruídos e melhorar a qualidade dos dados, foi adotada uma nova estratégia baseada em pontos vetoriais gerados manualmente no QGIS.

Foram utilizados dois arquivos shapefile:

\begin{itemize}
    \item \texttt{pontos\_mistos.shp}: contendo amostras das classes 0 e 1;
    \item \texttt{pontos\_nao\_qmd.shp}: contendo apenas amostras da classe 0.
\end{itemize}

A identificação das classes foi realizada a partir da coluna:

\begin{center}
\texttt{class}
\end{center}

Durante o pré-processamento:

\begin{itemize}
    \item Os pontos foram convertidos para coordenadas raster;
    \item Foram extraídos apenas os pixels correspondentes aos pontos amostrados;
    \item Pixels localizados em regiões \textit{NoData} foram removidos;
    \item As variáveis foram normalizadas antes do treinamento.
\end{itemize}

\subsection{Índices espectrais utilizados}

Foram utilizadas imagens Landsat pré e pós-fogo:

\begin{itemize}
    \item Pré-fogo: Landsat 8;
    \item Pós-fogo: Landsat 9.
\end{itemize}

A partir das bandas espectrais foram calculados os seguintes índices:

\begin{itemize}
    \item \textbf{dNBR} --- Delta Normalized Burn Ratio;
    \item \textbf{dNBRSWIR};
    \item \textbf{dNDVI} --- Delta Normalized Difference Vegetation Index.
\end{itemize}

Esses índices representam alterações espectrais relacionadas à perda de vegetação e à severidade da queima.

\subsection{Cena utilizada no treinamento}

A Figura~\ref{fig:cena_treinamento} apresenta a cena Landsat utilizada no treinamento do modelo KNN.

\begin{figure}[H]
    \centering
    \includegraphics[width=0.95\linewidth]{cena_treinamento.png}
    \caption{Cena Landsat utilizada no treinamento do modelo KNN.}
    \label{fig:cena_treinamento}
\end{figure}

\subsection{Resultados do treinamento}

O conjunto de dados foi dividido em:

\begin{itemize}
    \item 70\% para treinamento;
    \item 30\% para teste.
\end{itemize}

Os resultados obtidos durante a validação interna foram:

\begin{table}[H]
\centering
\caption{Métricas obtidas no conjunto de treinamento/teste interno.}
\begin{tabular}{lccc}
\hline
Classe & Precision & Recall & F1-score \\
\hline
0 (Não queimada) & 0.76 & 0.69 & 0.72 \\
1 (Queimada)     & 0.71 & 0.77 & 0.74 \\
\hline
\end{tabular}
\end{table}

O modelo apresentou:

\begin{itemize}
    \item Accuracy: 0.73;
    \item F1-score médio: 0.73.
\end{itemize}

Esses resultados indicaram desempenho moderadamente satisfatório dentro da própria cena utilizada no treinamento.

\subsection{Avaliação em nova cena}

Após o treinamento, o modelo foi aplicado em uma nova cena Landsat, distinta daquela utilizada durante o treinamento.

A Figura~\ref{fig:cena_teste} apresenta a cena de teste utilizada para avaliar a capacidade de generalização espacial do modelo.

\begin{figure}[H]
    \centering
    \includegraphics[width=0.95\linewidth]{cena_teste.png}
    \caption{Cena Landsat utilizada para teste do modelo KNN.}
    \label{fig:cena_teste}
\end{figure}

A Figura~\ref{fig:cena_classificada} apresenta o resultado da classificação produzida pelo modelo KNN na nova cena.

\begin{figure}[H]
    \centering
    \includegraphics[width=0.95\linewidth]{cena_teste_classificada.png}
    \caption{Resultado da classificação da nova cena utilizando o modelo KNN.}
    \label{fig:cena_classificada}
\end{figure}

Os resultados quantitativos obtidos foram:

\begin{table}[H]
\centering
\caption{Métricas obtidas na nova cena Landsat.}
\begin{tabular}{lccc}
\hline
Classe & Precision & Recall & F1-score \\
\hline
0 (Não queimada) & 0.99 & 0.69 & 0.81 \\
1 (Queimada)     & 0.09 & 0.82 & 0.17 \\
\hline
\end{tabular}
\end{table}

Métricas globais:

\begin{itemize}
    \item Accuracy: 0.69;
    \item Precision: 0.094;
    \item Recall: 0.825;
    \item F1-score: 0.169.
\end{itemize}

\subsection{Análise dos resultados}

Embora o modelo tenha apresentado desempenho razoável durante o treinamento, os resultados obtidos na nova cena foram significativamente inferiores.

O principal problema observado foi a baixa capacidade de generalização espacial do modelo. Na prática, verificou-se:

\begin{itemize}
    \item Grande quantidade de falsos positivos;
    \item Superestimação das áreas queimadas;
    \item Sensibilidade às diferenças radiométricas entre cenas;
    \item Forte dependência das características específicas da cena de treinamento.
\end{itemize}

A precisão da classe ``Queimada'' caiu para aproximadamente 9\%, indicando que a maior parte dos pixels classificados como queimados eram, na realidade, áreas não queimadas.

Apesar do \textit{recall} elevado (0.82), o modelo apresentou baixa confiabilidade prática devido à elevada taxa de falsos positivos.

\subsection{Comparação entre treinamento e teste}

A Figura~\ref{fig:comparacao_metricas} apresenta uma comparação entre as principais métricas obtidas no treinamento e na nova cena de teste.

\begin{figure}[H]
    \centering
    \includegraphics[width=0.9\linewidth]{grafico_metricas_knn.png}
    \caption{Comparação das métricas entre treinamento e teste em nova cena.}
    \label{fig:comparacao_metricas}
\end{figure}

Observa-se que, embora a acurácia geral tenha permanecido relativamente próxima, houve uma queda significativa na precisão e no F1-score da classe ``Queimada'', evidenciando a limitação do modelo em cenários distintos do treinamento.

\subsection{Conclusão}

Os resultados demonstraram que o modelo KNN conseguiu aprender parcialmente os padrões espectrais presentes na cena de treinamento, alcançando aproximadamente 73\% de acurácia.

Entretanto, quando aplicado em novas cenas Landsat, o modelo apresentou forte degradação de desempenho, principalmente devido ao aumento expressivo de falsos positivos.

Esses resultados indicam que o KNN apresentou baixa robustez para generalização espacial em aplicações de sensoriamento remoto para detecção de áreas queimadas.


\end{document}